\setcounter{section}{2}

\Needspace{5\baselineskip}
\hypertarget{offline-reinforcement-learning}{%
\section{Offline Reinforcement Learning}\label{offline-reinforcement-learning}}

\Needspace{5\baselineskip}
\hypertarget{i.-definition-and-main-methods-of-offline-reinforcement-learning}{%
\subsection{I. Definition and Main Methods of Offline Reinforcement Learning}\label{i.-definition-and-main-methods-of-offline-reinforcement-learning}}

The core of offline reinforcement learning is that the agent learns a policy only from a previously collected static dataset in an experience replay buffer, without interacting with the environment. This addresses excessively costly online trial and error or delayed feedback in real-world settings (such as autonomous driving, healthcare, and recommender systems).

In on-policy or off-policy algorithms, the agent can correct errors through interaction. In a purely offline setting, however, directly using a conventional off-policy algorithm (such as DDPG) fails because the distribution of state-action pairs \((s,a)\) in the dataset does not match the distribution generated by the current policy \(\pi\). The value function \(Q(s,a)\) assigns incorrectly high values to actions absent from the dataset (out-of-distribution, OOD). Because the agent cannot verify them in the environment, these errors continue accumulating and cause the policy to fail.

There are two main approaches to this problem: constrain policy choices (as in BCQ) or constrain the value function (as in CQL).

\Needspace{5\baselineskip}
\hypertarget{ii.-batch-constrained-q-learning-bcq}{%
\subsection{II. Batch-Constrained Q-Learning (BCQ)}\label{ii.-batch-constrained-q-learning-bcq}}

BCQ's core idea is to address ``extrapolation error'' by constraining the policy to the distribution of the offline dataset, or batch, and finding the best action within that range.

\Needspace{5\baselineskip}
\hypertarget{i-model-components}{%
\subsubsection{(I) Model Components}\label{i-model-components}}

\Needspace{5\baselineskip}
BCQ mainly consists of three types of networks:

\begin{enumerate}
\def\labelenumi{\arabic{enumi}.}
\tightlist
\item
  Critic Q network: estimates action value \(Q_\theta(s,a)\). Implementations usually use two Q networks, \(Q_{\theta_1}(s,a)\) and \(Q_{\theta_2}(s,a)\), combined conservatively to mitigate Q-value overestimation.
\item
  Conditional VAE: learns the distribution of possible actions in state \(s\) in the offline dataset, approximating behavior policy \(P_{\mathcal B}(a \mid s)\). It consists of encoder \(E_\omega(s,a)\) and decoder \(D_\omega(s,z)\).
\item
  Perturbation network: denoted \(\xi_\phi(s,a)\). It takes candidate actions generated by the VAE and adjusts them within a very small range, keeping them near the data distribution while obtaining higher \(Q\) values.
\end{enumerate}

\Needspace{5\baselineskip}
\hypertarget{ii-workflow}{%
\subsubsection{(II) Workflow}\label{ii-workflow}}

\hypertarget{train-the-vae}{%
\paragraph{1. Train the VAE}\label{train-the-vae}}

The first step is to train the VAE, which forms the basis of BCQ's action-range constraint.

\begin{itemize}
\tightlist
\item
  Training data are real \((s,a)\) pairs from the offline buffer.
\item
  The VAE learns ``what common actions \(a\) in the dataset look like in state \(s\).''
\item
  The encoder generates a latent-variable distribution from \((s,a)\), and the decoder reconstructs action \(\hat{a}\) from \((s,z)\).
\item
  After training, given a new state \(s\) and different sampled \(z\), the decoder generates a batch of candidate actions close to the data distribution.
\end{itemize}

The specific training method is described later.

\hypertarget{calculate-the-target-q-value}{%
\paragraph{2. Calculate the Target Q Value}\label{calculate-the-target-q-value}}

In Q-learning, the target \(Q\) value is the maximum of ``reward + \(Q\) value at next state \(s'\).'' The difficulty here is estimating the \(Q\) value of the next state \(s'\).

\Needspace{5\baselineskip}
BCQ does not calculate \(\max_a Q(s',a)\) directly over all possible actions, because that can select OOD actions outside the offline dataset. Instead, it maximizes only over candidate actions generated by the VAE:

\begin{enumerate}
\def\labelenumi{\arabic{enumi}.}
\tightlist
\item
  For next state \(s'\), sample \(n\) latent variables \(z_i \sim \mathcal{N}(0,I)\) from the standard normal distribution.
\item
  Generate candidate actions \(a_i = D_{\omega'}(s',z_i)\) with the target VAE.
\item
\Needspace{5\baselineskip}
  Make small corrections with the target perturbation network:
\end{enumerate}

\[
\tilde a_i = a_i + \xi_{\phi'}(s',a_i)
\]

\begin{enumerate}
\def\labelenumi{\arabic{enumi}.}
\setcounter{enumi}{3}
\tightlist
\item
\Needspace{5\baselineskip}
  Select the candidate with the highest target Q value. A common BCQ target is:
\end{enumerate}

\[
y = r + \gamma \max_i \bigl[
\lambda \min_{j=1,2} Q_{\theta'_j}(s', \tilde a_i)
+ (1-\lambda) \max_{j=1,2} Q_{\theta'_j}(s', \tilde a_i)
\bigr]
\]

Here \(\lambda\) controls the target's conservatism, and \(\theta'_j\), \(\omega'\), and \(\phi'\) denote target-network parameters.

\hypertarget{update-the-q-network}{%
\paragraph{3. Update the Q Network}\label{update-the-q-network}}

\Needspace{5\baselineskip}
The Q network updates using the standard TD error, bringing the current \(Q\) value closer to target \(y\):

\[
L_Q(\theta)=\mathbb{E}_{(s,a,r,s')\sim\mathcal B}
[(Q_\theta(s,a)-y)^{2}]
\]

With twin \(Q\) networks, separately minimize the mean squared errors for \(Q_{\theta_1}\) and \(Q_{\theta_2}\).

\hypertarget{update-the-perturbation-network-see-below}{%
\paragraph{4. Update the Perturbation Network (See Below)}\label{update-the-perturbation-network-see-below}}

\Needspace{5\baselineskip}
\hypertarget{iii-training-the-vae}{%
\subsubsection{(III) Training the VAE}\label{iii-training-the-vae}}

\Needspace{5\baselineskip}
VAE training involves four steps: encoding, reparameterization, decoding, and loss optimization:

\begin{enumerate}
\def\labelenumi{\arabic{enumi}.}
\tightlist
\item
  The encoder takes \((s,a)\) as input and outputs latent-distribution parameters \(\mu_\omega(s,a)\) and \(\sigma_\omega(s,a)\).
\item
\Needspace{5\baselineskip}
  Sample a latent variable through reparameterization:
\end{enumerate}

\[
\epsilon \sim \mathcal{N}(0,I), \qquad z=\mu_\omega(s,a)+\sigma_\omega(s,a)\odot\epsilon
\]

\Needspace{4\baselineskip}
\begin{enumerate}
\def\labelenumi{\arabic{enumi}.}
\setcounter{enumi}{2}
\tightlist
\item
\Needspace{5\baselineskip}
  The decoder takes \((s,z)\) as input and outputs the reconstructed action:
\end{enumerate}

\[
\hat a = D_\omega(s,z)
\]

\Needspace{4\baselineskip}
\begin{enumerate}
\def\labelenumi{\arabic{enumi}.}
\setcounter{enumi}{3}
\tightlist
\item
\Needspace{5\baselineskip}
  The loss consists of action reconstruction error and a KL-divergence regularizer:
\end{enumerate}

\[
L_{\mathrm{VAE}}
=\mathbb{E}[\lVert a-D_\omega(s,z)\rVert^{2}]
+ D_{\mathrm{KL}}(q_\omega(z\mid s,a)\Vert \mathcal{N}(0,I))
\]

Reconstruction error enables the decoder to generate actions appearing in the dataset. KL regularization brings the latent space close to the standard normal distribution, allowing direct sampling from \(\mathcal{N}(0,I)\) at test time.

\Needspace{5\baselineskip}
\hypertarget{iv-why-use-a-vae-instead-of-regression-or-cross-entropy}{%
\subsubsection{(IV) Why Use a VAE Instead of Regression or Cross-Entropy?}\label{iv-why-use-a-vae-instead-of-regression-or-cross-entropy}}

\hypertarget{the-problem-with-regression}{%
\paragraph{1. The Problem with Regression}\label{the-problem-with-regression}}

An ordinary regression network usually learns a deterministic mapping \(a=\pi(s)\). If multiple reasonable actions exist in the same state \(s\), regression tends to output their average, which may not be a reasonable action.

For example, in intersection navigation, historical data for the same state \(s\) may include turning left \(a_1=-1\) and turning right \(a_2=+1\), both of which avoid obstacles. If a deterministic network \(\mathrm{Net}(s)\) minimizes mean squared error:

\[
Loss=(Net(s)-(-1))^{2}+(Net(s)-(+1))^{2}
\]

The optimal output is close to \(a=0\), meaning going straight. Going straight may be precisely the dangerous action. This illustrates that the mean of a multimodal action distribution can fall in a low-probability or even incorrect region.

\hypertarget{the-problem-with-cross-entropy}{%
\paragraph{2. The Problem with Cross-Entropy}\label{the-problem-with-cross-entropy}}

Cross-entropy suits discrete-action classification, but continuous control actions cannot be directly represented by a finite set of classes. Forced discretization quickly causes the number of classes to explode.

\Needspace{5\baselineskip}
For example, suppose each joint angle of a robotic arm lies in \([-180^{\circ},180^{\circ}]\). Dividing each joint into 100 bins gives the following number of joint classes for 7 joints:

\[
100^{7} = 100000000000000
\]

The Softmax output layer would have to cover an enormous combinatorial action space, making computation and generalization difficult. A VAE is therefore more suitable than simple regression or cross-entropy classification for continuous, multimodal action distributions.

\hypertarget{benefits-of-a-vae}{%
\paragraph{3. Benefits of a VAE}\label{benefits-of-a-vae}}

Together, the VAE's encoder and decoder effectively form an ``imitator'' that can efficiently ``imitate'' a batch of similar \(a\) given \(a\), regardless of the distribution of \(a\).

A VAE uses latent variable \(z\) to represent multiple possible actions in the same state. In BCQ's conditional VAE, the mapping is not simply \(s\to a\) but \((s,z)\to a\).

\begin{itemize}
\tightlist
\item
  During training, the encoder sees the complete \((s,a)\). If the same state has both ``turn left'' and ``turn right'' actions, it can map ``state + left turn'' to some \(z_{\mathrm{left}}\) and ``state + right turn'' to another \(z_{\mathrm{right}}\).
\item
  During generation, only state \(s\) is given; different \(z\) are sampled from \(\mathcal{N}(0,I)\) and passed to decoder \(D_\omega(s,z)\) to generate different actions.
\end{itemize}

The key remains that changing latent variable \(z\) lets the VAE generate multiple different actions in the same state that all remain close to the data distribution.

\Needspace{5\baselineskip}
Continuing the example, the decoder can map different latent variables back to different actions:

\begin{itemize}
\tightlist
\item
  When \(z_1\) is sampled, \(D_\omega(s,z_1)\) may output \(a\approx -1\), corresponding to turning left.
\item
  When \(z_2\) is sampled, \(D_\omega(s,z_2)\) may output \(a\approx +1\), corresponding to turning right.
\end{itemize}

Thus, the VAE learns conditional action distribution \(P(a\mid s)\) instead of forcing state \(s\) into a single fixed action. By changing \(z\), BCQ generates multiple candidate actions in the same state that are ``seen in the data or close to those seen.''

\Needspace{5\baselineskip}
\hypertarget{iv-updating-the-perturbation-network}{%
\subsubsection{(IV) Updating the Perturbation Network}\label{iv-updating-the-perturbation-network}}

The VAE keeps actions near the data distribution, but it only imitates behavior data and does not guarantee that generated actions are optimal under the current Q function. The perturbation network makes local improvements under the constraint of ``not deviating too far from the data distribution.''

\begin{itemize}
\tightlist
\item
  Input: VAE-generated action \(a_{\mathrm{vae}}=D_\omega(s,z)\).
\item
  Output: a small perturbation \(\xi_\phi(s,a_{\mathrm{vae}})\).
\item
\Needspace{5\baselineskip}
  Corrected action:
\end{itemize}

\[
\tilde a = a_{\mathrm{vae}}+\xi_\phi(s,a_{\mathrm{vae}})
\]

\begin{itemize}
\tightlist
\item
  Constraint: the perturbation is usually limited to \([-\Phi,\Phi]\), such as \(\Phi=0.05\), embodying the batch-constrained idea.
\end{itemize}

\Needspace{5\baselineskip}
The perturbation network usually updates after the Q network or alternately with it. During its update, freeze the Q network and optimize only perturbation parameters \(\phi\), aiming to obtain higher Q values for the adjusted actions:

\[
L_\phi
=-\mathbb{E}_{s\sim\mathcal B,\ z\sim\mathcal{N}(0,I)}
\bigl[
Q_\theta(s,a_{\mathrm{vae}}+\xi_\phi(s,a_{\mathrm{vae}}))
\bigr]
\]

Here \(a_{\mathrm{vae}}=D_\omega(s,z)\). Minimizing \(L_\phi\) is equivalent to maximizing the adjusted action's \(Q\) value.

\Needspace{5\baselineskip}
\hypertarget{iii.-conservative-q-learning-cql}{%
\subsection{III. Conservative Q-Learning (CQL)}\label{iii.-conservative-q-learning-cql}}

CQL builds on SAC by modifying how the Q function updates. Since unknown actions are easily overestimated when interaction is unavailable, it adds a penalty directly to the Q function's update objective, artificially lowering Q values for unseen actions.

\Needspace{5\baselineskip}
CQL's Critic loss can be understood as ``Bellman error + conservative regularization'':

\[
L_{\mathrm{CQL}}(\theta)
=\frac{1}{2}\mathbb{E}_{(s,a,r,s')\sim\mathcal D}
[(Q_\theta(s,a)-y)^{2}]
+\alpha \mathcal R(\theta)
\]

The first term is standard Bellman error, making the Q function satisfy the temporal-difference target. The second term, \(\mathcal R(\theta)\), is a conservative regularizer that lowers potentially overestimated actions, particularly OOD actions rarely appearing in the offline dataset.

\Needspace{5\baselineskip}
For discrete action spaces, the conservative regularizer is commonly written as:

\[
\mathcal R(\theta)
=\mathbb{E}_{s\sim\mathcal D}
\bigl[
\log \sum_a \exp(Q_\theta(s,a))
-\mathbb{E}_{a\sim\mathcal D(\cdot\mid s)}[Q_\theta(s,a)]
\bigr]
\]

\Needspace{5\baselineskip}
This contains two opposing forces:

\begin{enumerate}
\def\labelenumi{\arabic{enumi}.}
\tightlist
\item
  \(-\mathbb{E}_{a\sim\mathcal D(\cdot\mid s)}[Q_\theta(s,a)]\): minimizing the loss raises the Q values of actual actions in the dataset.
\item
  \(\log \sum_a \exp(Q_\theta(s,a))\): this is LogSumExp, also viewed as a Soft-Maximum, which approximately focuses on actions with the highest Q values in the current state. Minimizing it lowers high-Q actions, especially OOD actions insufficiently supported by the dataset.
\end{enumerate}

The final effect is that actual dataset actions are supported by the data term and are not indiscriminately suppressed; OOD actions lack that support and are mainly suppressed by LogSumExp, producing a more conservative Q function.

\Needspace{5\baselineskip}
In continuous action spaces, the sum or integral corresponding to \(\log\sum_a\exp(Q(s,a))\) cannot be calculated exactly and directly, so CQL usually approximates it through sampling. An implementation-level expression is:

\[
L(\theta)
=\alpha\mathbb{E}_{s\sim\mathcal D}
\bigl[
\mathbb{E}_{a\sim\mu(a\mid s)}[Q_\theta(s,a)]
-\mathbb{E}_{a\sim\mathcal D(\cdot\mid s)}[Q_\theta(s,a)]
\bigr]
+\mathrm{Bellman\ Error}
\]

\Needspace{5\baselineskip}
Here \(\mu(a\mid s)\) is the sampling distribution for generating ``potentially overestimated actions.'' In practice, two types of actions are commonly sampled to approximate the conservative term:

\begin{itemize}
\tightlist
\item
  Current-policy actions: \(a\sim\pi_\phi(\cdot\mid s)\), to lower actions that the policy may choose but that lack data support.
\item
  Random actions: \(a\sim\mathrm{Uniform}(-1,1)\), to cover a broader continuous action space.
\end{itemize}

Thus, CQL does not constrain policy actions directly. It trains a more conservative Q function so that policy optimization naturally avoids overestimated OOD actions.

\FloatBarrier
\Needspace{5\baselineskip}
\hypertarget{references-64}{%
\subsection{References}\label{references-64}}

\vspace{0.25em}
\begin{itemize}
\tightlist
\item
  \href{https://hrl.boyuai.com/}{Hands-on Reinforcement Learning (translated title; in Chinese)}.
\item
  Fujimoto, S., Meger, D., \& Precup, D. (2019). \href{https://arxiv.org/abs/1812.02900}{Off-Policy Deep Reinforcement Learning without Exploration}. ICML 2019.
\item
  Kumar, A., Zhou, A., Tucker, G., \& Levine, S. (2020). \href{https://arxiv.org/abs/2006.04779}{Conservative Q-Learning for Offline Reinforcement Learning}. NeurIPS 2020.
\end{itemize}

\clearpage
